%%%%%%%%%%%%%%%%
%%% exod 12 %%%
%%%%%%%%%%%%%%%%

\Chapter{12}

\Verse{1}
{
	\hP{וַיֹּאמֶר יְהֹוָה אֶל מֹשֶׁה וְאֶל אַהֲרֹן בְּאֶרֶץ מִצְרַיִם לֵאמֹר׃}
}
{
	\eP{
		12 And YHWH said to Moses and to Aaron in the land of
		Egypt, saying,
	}
}
{
}

\Verse{2}
{
	\hP{הַחֹדֶשׁ הַזֶּה לָכֶם רֹאשׁ חֳדָשִׁים רִאשׁוֹן הוּא לָכֶם לְחדְשֵׁי הַשָּׁנָה׃}
}
{
	\eP{
		“This month is the beginning of months for you. It is
		first of the months of the year for you.
	}
}
{
%	\footnote{
%		first of the months of the year. According to the Torah, the new year
%		begins in spring, not in the fall as Jews celebrate it now. The holiday
%		that comes on the first day of fall is never referred to in the Torah as
%		Rosh Hashanah (Lev 23:23–25). 12:2–11. This month … This is sometimes
%		thought to be the instruction for the annual observance of Passover for
%		all time. But the wording and context seem rather to apply only to the
%		night before the actual exodus from Egypt. The Passover instructions for
%		future generations come later, starting at 12:14.
%	}
}

\Verse{3}
{
	\hP{דַּבְּרוּ אֶל כּל עֲדַת יִשְׂרָאֵל לֵאמֹר בֶּעָשֹׂר לַחֹדֶשׁ הַזֶּה וְיִקְחוּ לָהֶם אִישׁ שֶׂה לְבֵית אָבֹת שֶׂה לַבָּיִת׃}
}
{
	\eP{
		Speak to all of the congregation of Israel, saying: On the
		tenth of this month, let them each take a lamb for the
		fathers’ houses, a lamb per house.
	}
}
{
}

\Verse{4}
{
	\hP{וְאִם יִמְעַט הַבַּיִת מִהְיוֹת מִשֶּׂה וְלָקַח הוּא וּשְׁכֵנוֹ הַקָּרֹב אֶל בֵּיתוֹ בְּמִכְסַת נְפָשֹׁת אִישׁ לְפִי אכְלוֹ תָּכֹסּוּ עַל הַשֶּׂה׃}
}
{
	\eP{
		And if the household will be too few for a lamb, then he
		and his neighbor who is close to his house will take it
		according to the count of persons; you shall count each
		person according to what he eats for the lamb.
	}
}
{
}

\Verse{5}
{
	\hP{שֶׂה תָמִים זָכָר בֶּן שָׁנָה יִהְיֶה לָכֶם מִן הַכְּבָשִׂים וּמִן הָעִזִּים תִּקָּחוּ׃}
}
{
	\eP{
		You shall have an unblemished, male, year-old lamb; you
		shall take it from the sheep or from the goats.
	}
}
{
}

\Verse{6}
{
	\hP{וְהָיָה לָכֶם לְמִשְׁמֶרֶת עַד אַרְבָּעָה עָשָׂר יוֹם לַחֹדֶשׁ הַזֶּה וְשָׁחֲטוּ אֹתוֹ כֹּל קְהַל עֲדַת יִשְׂרָאֵל בֵּין הָעַרְבָּיִם׃}
}
{
	\eP{
		And it will be for you to watch over until the fourteenth
		day of this month. And all the community of the
		congregation of Israel will slaughter it ‘between the two
		evenings.’
	}
}
{
%	\footnote{
%		between the two evenings. Understood to mean the early evening hours:
%		“dusk” or “twilight.”
%	}
}

\Verse{7}
{
	\hP{וְלָקְחוּ מִן הַדָּם וְנָתְנוּ עַל שְׁתֵּי הַמְּזוּזֹת וְעַל הַמַּשְׁקוֹף עַל הַבָּתִּים אֲשֶׁר יֹאכְלוּ אֹתוֹ בָּהֶם׃}
}
{
	\eP{
		And they will take some of the blood and place it on the
		two doorposts and on the lintel on the houses in which
		they will eat it.
	}
}
{
}

\Verse{8}
{
	\hP{וְאָכְלוּ אֶת הַבָּשָׂר בַּלַּיְלָה הַזֶּה צְלִי אֵשׁ וּמַצּוֹת עַל מְרֹרִים יֹאכְלֻהוּ׃}
}
{
	\eP{
		And they will eat the meat in this night; they will eat it
		fire-roasted and with unleavened bread on bitter herbs.
	}
}
{
}

\Verse{9}
{
	\hP{אַל תֹּאכְלוּ מִמֶּנּוּ נָא וּבָשֵׁל מְבֻשָּׁל בַּמָּיִם כִּי אִם צְלִי אֵשׁ רֹאשׁוֹ עַל כְּרָעָיו וְעַל קִרְבּוֹ׃}
}
{
	\eP{
		Do not eat any of it raw or cooked in water, but
		fire-roasted: its head with its legs and with its innards.
	}
}
{
}

\Verse{10}
{
	\hP{וְלֹא תוֹתִירוּ מִמֶּנּוּ עַד בֹּקֶר וְהַנֹּתָר מִמֶּנּוּ עַד בֹּקֶר בָּאֵשׁ תִּשְׂרֹפוּ׃}
}
{
	\eP{
		And do not leave any of it until morning; and you shall
		burn what is left of it until morning in fire.
	}
}
{
}

\Verse{11}
{
	\hP{וְכָכָה תֹּאכְלוּ אֹתוֹ מתְנֵיכֶם חֲגֻרִים נַעֲלֵיכֶם בְּרַגְלֵיכֶם וּמַקֶּלְכֶם בְּיֶדְכֶם וַאֲכַלְתֶּם אֹתוֹ בְּחִפָּזוֹן פֶּסַח הוּא לַיהֹוָה׃}
}
{
	\eP{
		And you shall eat it like this: your hips clothed, your
		shoes on your feet, and your staff in your hand; and you
		shall eat it in haste. It is YHWH’s Passover.
	}
}
{
}

\Verse{12}
{
	\hP{וְעָבַרְתִּי בְאֶרֶץ מִצְרַיִם בַּלַּיְלָה הַזֶּה וְהִכֵּיתִי כל בְּכוֹר בְּאֶרֶץ מִצְרַיִם מֵאָדָם וְעַד בְּהֵמָה וּבְכל אֱלֹהֵי מִצְרַיִם אֶעֱשֶׂה שְׁפָטִים אֲנִי יְהֹוָה׃}
}
{
	\eP{
		“And I shall pass through the land of Egypt in this night,
		and I shall strike every firstborn in the land of Egypt,
		from human to animal, and I shall make judgments on all
		the gods of Egypt. I am YHWH.
	}
}
{
%	\footnote{
%		judgments on all the gods of Egypt. The best-known deity of Egyptian
%		religion is the sun, and Egyptian religion was profoundly concerned with
%		death. The plagues culminate with “judgments” on these. In the ninth
%		plague there is darkness for three days. And in the final plague death
%		itself is shown to be in YHWH’s control, as only firstborn humans and
%		animals die. (The first plague turns Egypt’s waters to blood, and this
%		too undermines Egyptian deities.) Some readers may see in this an
%		implicit recognition of the existence of the pagan deities, who must be
%		thought to exist in order to be thus defeated. The text, however, does
%		not present these “judgments” as a defeat of the gods. The plagues are
%		rather a show of where power resides, namely, outside the gods (which is
%		to say, outside of nature), beyond them. The forces of nature are not
%		personified in the plagues narrative, and they do not confront or
%		challenge the God of Israel. They are merely manipulated in the course
%		of events.
%	}
}

\Verse{13}
{
	\hP{וְהָיָה הַדָּם לָכֶם לְאֹת עַל הַבָּתִּים אֲשֶׁר אַתֶּם שָׁם וְרָאִיתִי אֶת הַדָּם וּפָסַחְתִּי עֲלֵכֶם וְלֹא יִהְיֶה בָכֶם נֶגֶף לְמַשְׁחִית בְּהַכֹּתִי בְּאֶרֶץ מִצְרָיִם׃}
}
{
	\eP{
		And the blood will be as a sign for you on the houses in
		which you are, and I shall see the blood, and I shall halt
		at you, and there won’t be a plague among you as a
		destroyer when I strike in the land of Egypt. 12:13,23,27.
		halt. Hebrew psa does not mean “to pass over.” That
		wording has led people to images of the deity floating
		over houses. The verb means “to halt” or “to walk in a
		halting manner”; it can refer to limping (2 Sam 4:4). The
		noun form pissa means a cripple (2 Sam 9:13). Admittedly,
		this verb occurs in Isaiah in a verse that pictures {\God}
		defending Jerusalem “like birds flying” (31:5). Still,
		“halting” fits with the context here in Exodus, especially
		in 12:23, where it suggests a conception of the deity
		moving along through Egypt, spotting the blood on the
		doorposts, and coming to an abrupt stop. {\God} “halts on the
		threshold” and does not allow the destroying force to
		enter the house. “Passing over” the threshold does not
		really fit with this picture of blocking or preventing the
		destroyer. 12:13,23. destroyer. The Hebrew pun conveys the
		message: the Israelites are to slaughter (root ) the lamb
		so as to keep out the destroyer (root ).
	}
}
{
}

\Verse{14}
{
	\hP{וְהָיָה הַיּוֹם הַזֶּה לָכֶם לְזִכָּרוֹן וְחַגֹּתֶם אֹתוֹ חַג לַיהֹוָה לְדֹרֹתֵיכֶם חֻקַּת עוֹלָם תְּחגֻּהוּ׃}
}
{
	\eP{
		“And this day will become a commemoration for you, and you
		shall celebrate it, a festival to YHWH; you shall
		celebrate it through your generations, an eternal law:
	}
}
{
}

\Verse{15}
{
	\hP{שִׁבְעַת יָמִים מַצּוֹת תֹּאכֵלוּ אַךְ בַּיּוֹם הָרִאשׁוֹן תַּשְׁבִּיתוּ שְּׂאֹר מִבָּתֵּיכֶם כִּי כּל אֹכֵל חָמֵץ וְנִכְרְתָה הַנֶּפֶשׁ הַהִוא מִיִּשְׂרָאֵל מִיּוֹם הָרִאשֹׁן עַד יוֹם הַשְּׁבִעִי׃}
}
{
	\eP{
		Seven days you shall eat unleavened bread. Indeed, on the
		first day you shall make leaven cease from your houses.
		Because anyone who eats leavened bread: that person will
		be cut off from Israel—from the first day to the seventh
		day.
	}
}
{
%	\footnote{
%		make leaven cease. In the preceding verse it is said that the Passover
%		festival is to be celebrated as a commemoration of what happened in
%		Egypt. The wording now conveys this: The command to “make leaven cease”
%		recalls Pharaoh’s first meeting with Moses, in which he says, “You’ve
%		made them cease from their burdens” (Exod 5:5).
%	}
}

\Verse{16}
{
	\hP{וּבַיּוֹם הָרִאשׁוֹן מִקְרָא קֹדֶשׁ וּבַיּוֹם הַשְּׁבִיעִי מִקְרָא קֹדֶשׁ יִהְיֶה לָכֶם כּל מְלָאכָה לֹא יֵעָשֶׂה בָהֶם אַךְ אֲשֶׁר יֵאָכֵל לְכל נֶפֶשׁ הוּא לְבַדּוֹ יֵעָשֶׂה לָכֶם׃}
}
{
	\eP{
		And you will have a holy assembly on the first day and a
		holy assembly on the seventh day. Not any work will be
		done on them. Just what will be eaten by each person: that
		alone will be done for you.
	}
}
{
}

\Verse{17}
{
	\hP{וּשְׁמַרְתֶּם אֶת הַמַּצּוֹת כִּי בְּעֶצֶם הַיּוֹם הַזֶּה הוֹצֵאתִי אֶת צִבְאוֹתֵיכֶם מֵאֶרֶץ מִצְרָיִם וּשְׁמַרְתֶּם אֶת הַיּוֹם הַזֶּה לְדֹרֹתֵיכֶם חֻקַּת עוֹלָם׃}
}
{
	\eP{
		And you shall observe the unleavened bread, because in
		this very day I brought out your masses from the land of
		Egypt, and you shall observe this day through your
		generations, an eternal law.
	}
}
{
%	\footnote{
%		observe the unleavened bread. Meaning: observe the commandment
%		concerning the unleavened bread. The Septuagint and the Samaritan have
%		“the commandment.” The confusion may owe to the fact that the words
%		“unleavened bread” and “commandment” look the same in Hebrew consonants:
%		.
%	}
}

\Verse{18}
{
	\hP{בָּרִאשֹׁן בְּאַרְבָּעָה עָשָׂר יוֹם לַחֹדֶשׁ בָּעֶרֶב תֹּאכְלוּ מַצֹּת עַד יוֹם הָאֶחָד וְעֶשְׂרִים לַחֹדֶשׁ בָּעָרֶב׃}
}
{
	\eP{
		In the first month, on the fourteenth day of the month, in
		the evening, you shall eat unleavened bread, until the
		twenty-first day of the month, in the evening.
	}
}
{
}

\Verse{19}
{
	\hP{שִׁבְעַת יָמִים שְׂאֹר לֹא יִמָּצֵא בְּבָתֵּיכֶם כִּי כּל אֹכֵל מַחְמֶצֶת וְנִכְרְתָה הַנֶּפֶשׁ הַהִוא מֵעֲדַת יִשְׂרָאֵל בַּגֵּר וּבְאֶזְרַח הָאָרֶץ׃}
}
{
	\eP{
		Seven days leaven shall not be found in your houses,
		because anyone who eats something leavened: that person
		will be cut off from the congregation of Israel, whether
		the alien or the citizen of the land.
	}
}
{
}

\Verse{20}
{
	\hP{כּל מַחְמֶצֶת לֹא תֹאכֵלוּ בְּכֹל מוֹשְׁבֹתֵיכֶם תֹּאכְלוּ מַצּוֹת׃}
}
{
	\eP{
		You shall not eat anything leavened; in all your homes you
		shall eat unleavened bread.”
	}
}
{
}

\Verse{21}
{
	\hJ{וַיִּקְרָא מֹשֶׁה לְכל זִקְנֵי יִשְׂרָאֵל וַיֹּאמֶר אֲלֵהֶם מִשְׁכוּ וּקְחוּ לָכֶם צֹאן לְמִשְׁפְּחֹתֵיכֶם וְשַׁחֲטוּ הַפָּסַח׃}
}
{
	\eJ{
		And Moses called all Israel’s elders, and he said to them,
		“Pull out and take a sheep for your families and slaughter
		the Passover.
	}
}
{
%	\footnote{
%		Passover. I have retained the English translation “Passover” rather than
%		using some form of the word “halt.” Although the latter would preserve
%		the connection between this word and the story of God’s halting, the
%		fact remains that Passover is now the established, famous name of the
%		holiday, and there is no point in ignoring this in a translation.
%	}
}

\Verse{22}
{
	\hJ{וּלְקַחְתֶּם אֲגֻדַּת אֵזוֹב וּטְבַלְתֶּם בַּדָּם אֲשֶׁר בַּסַּף וְהִגַּעְתֶּם אֶל הַמַּשְׁקוֹף וְאֶל שְׁתֵּי הַמְּזוּזֹת מִן הַדָּם אֲשֶׁר בַּסָּף וְאַתֶּם לֹא תֵצְאוּ אִישׁ מִפֶּתַח בֵּיתוֹ עַד בֹּקֶר׃}
}
{
	\eJ{
		And you’ll take a bunch of hyssop and dip in the blood
		that is in a basin and touch some of the blood that is in
		the basin to the lintel and to the two doorposts. And you
		shall not go out, each one, from his house’s entrance
		until morning.
	}
}
{
}

\Verse{23}
{
	\hJ{וְעָבַר יְהֹוָה לִנְגֹּף אֶת מִצְרַיִם וְרָאָה אֶת הַדָּם עַל הַמַּשְׁקוֹף וְעַל שְׁתֵּי הַמְּזוּזֹת וּפָסַח יְהֹוָה עַל הַפֶּתַח וְלֹא יִתֵּן הַמַּשְׁחִית לָבֹא אֶל בָּתֵּיכֶם לִנְגֹּף׃}
}
{
	\eJ{
		And YHWH will pass to strike Egypt, and He’ll see the
		blood on the lintel and on the two doorposts, and YHWH
		will halt at the entrance and will not allow the destroyer
		to come to your houses to strike.
	}
}
{
}

\Verse{24}
{
	\hJ{וּשְׁמַרְתֶּם אֶת הַדָּבָר הַזֶּה לְחק לְךָ וּלְבָנֶיךָ עַד עוֹלָם׃}
}
{
	\eJ{
		And you shall observe this thing as a law for you and your
		children forever.
	}
}
{
}

\Verse{25}
{
	\hJ{וְהָיָה כִּי תָבֹאוּ אֶל הָאָרֶץ אֲשֶׁר יִתֵּן יְהֹוָה לָכֶם כַּאֲשֶׁר דִּבֵּר וּשְׁמַרְתֶּם אֶת הָעֲבֹדָה הַזֹּאת׃}
}
{
	\eJ{
		And it will be, when you will come to the land that YHWH
		will give to you as He has spoken, that you shall observe
		this service.
	}
}
{
}

\Verse{26}
{
	\hJ{וְהָיָה כִּי יֹאמְרוּ אֲלֵיכֶם בְּנֵיכֶם מָה הָעֲבֹדָה הַזֹּאת לָכֶם׃}
}
{
	\eJ{
		And it will be, when your children will say to you, ‘What
		is this service to you?’
	}
}
{
}

\Verse{27}
{
	\hJ{וַאֲמַרְתֶּם זֶבַח פֶּסַח הוּא לַיהֹוָה אֲשֶׁר פָּסַח עַל בָּתֵּי בְנֵי יִשְׂרָאֵל בְּמִצְרַיִם בְּנגְפּוֹ אֶת מִצְרַיִם וְאֶת בָּתֵּינוּ הִצִּיל וַיִּקֹּד הָעָם וַיִּשְׁתַּחֲווּ׃}
}
{
	\eJ{
		that you shall say, ‘It is the Passover sacrifice to YHWH,
		because He halted at the houses of the children of Israel
		in Egypt when He struck Egypt, and He saved our houses.’”
		And the people knelt and bowed.
	}
}
{
}

\Verse{28}
{
	\hR{וַיֵּלְכוּ וַיַּעֲשׂוּ בְּנֵי יִשְׂרָאֵל כַּאֲשֶׁר צִוָּה יְהֹוָה אֶת מֹשֶׁה וְאַהֲרֹן כֵּן עָשׂוּ׃}
}
{
	\eR{
		And the children of Israel went and did as YHWH had
		commanded Moses and Aaron. They did so.
	}
}
{
}

\Verse{29}
{
	\hRJE{וַיְהִי בַּחֲצִי הַלַּיְלָה וַיהֹוָה הִכָּה כל בְּכוֹר בְּאֶרֶץ מִצְרַיִם מִבְּכֹר פַּרְעֹה הַיֹּשֵׁב עַל כִּסְאוֹ עַד בְּכוֹר הַשְּׁבִי אֲשֶׁר בְּבֵית הַבּוֹר וְכֹל בְּכוֹר בְּהֵמָה׃}
}
{
	\eRJE{
		And it was in the middle of the night, and YHWH struck
		every firstborn in the land of Egypt, from the firstborn
		of Pharaoh who was sitting on his throne to the firstborn
		of the prisoner who was in the prison house and every
		firstborn of an animal.
	}
}
{
}

\Verse{30}
{
	\hRJE{וַיָּקם פַּרְעֹה לַיְלָה הוּא וְכל עֲבָדָיו וְכל מִצְרַיִם וַתְּהִי צְעָקָה גְדֹלָה בְּמִצְרָיִם כִּי אֵין בַּיִת אֲשֶׁר אֵין שָׁם מֵת׃}
}
{
	\eRJE{
		And Pharaoh got up at night, he and all his servants and
		all Egypt, and there was a big cry in Egypt, because there
		was not a house in which there was not one dead.
	}
}
{
}

\Verse{31}
{
	\hRJE{וַיִּקְרָא לְמֹשֶׁה וּלְאַהֲרֹן לַיְלָה וַיֹּאמֶר קוּמוּ צְּאוּ מִתּוֹךְ עַמִּי גַּם אַתֶּם גַּם בְּנֵי יִשְׂרָאֵל וּלְכוּ עִבְדוּ אֶת יְהֹוָה כְּדַבֶּרְכֶם׃}
}
{
	\eRJE{
		And he called Moses and Aaron at night, and he said, “Get
		up. Go out from among my people, both you and the children
		of Israel, and go, serve YHWH as you spoke.
	}
}
{
%	\footnote{
%		he called Moses and Aaron. It is not clear whether this means that they
%		go back to him or that they just receive a message from him. The issue
%		is that Pharaoh had said, “Don’t continue to see my face. Because in the
%		day you see my face you’ll die!” And Moses had answered, “So you’ve
%		spoken: I won’t continue to see your face anymore” (10:28–29). So if
%		Pharaoh sees Moses now, it means that his threat has proved empty; and
%		Moses’ answer to him must be cynical: “So you’ve spoken!” Or, if Pharaoh
%		and Moses do not see each other now, then Moses’ answer must be
%		understood to be ironic: “You’re right, I won’t see you again, because
%		we’ll be gone!” Footnote 12:31. as you spoke. What we have seen here is
%		a gradual overwhelming of the power of Egypt through twists of events
%		and personalities, stymieing the Pharaoh’s initial confidence and
%		leading him to surrender unconditionally. In the first meeting he is
%		unbending. In the second meeting and through the first two plagues, he
%		stands firm as his magicians perform the wonders as well. The magicians
%		cannot perform the third plague, and so the Pharaoh, no longer
%		commanding cosmic forces himself, negotiates on the fourth (“Only don’t
%		go too far”). He gets away with his reneging on the negotiation, and so
%		he stands firm again through the fifth plague. His magicians themselves
%		suffer the sixth plague, and, following this indication that the
%		opposing power can even encroach on his forces, the Pharaoh talks to
%		Moses again on the seventh, generously sharing the blame with his
%		subjects (“I and my people are wrong …”). When he reneges yet again, and
%		an eighth plague strikes, his own servants urge him to give up resisting
%		(“Don’t you know yet that Egypt has perished?”). He then negotiates
%		again (“just the men”). When this only leads to a ninth plague,
%		darkening Egypt’s divine sun, he negotiates again, allowing the
%		children, but not the livestock, to go. When Moses insists on the
%		livestock as well, the Pharaoh returns to his original position,
%		refusing to let the people go and harshly telling Moses to get out and
%		never return. The horror of the tenth plague then comes, and the Pharaoh
%		capitulates utterly. It becomes clear that this had never been a matter
%		of negotiation at all, but rather an agonizing, gradual drawing of the
%		Pharaoh to a decision that had been inescapable from the start.
%	}
}

\Verse{32}
{
	\hRJE{גַּם צֹאנְכֶם גַּם בְּקַרְכֶם קְחוּ כַּאֲשֶׁר דִּבַּרְתֶּם וָלֵכוּ וּבֵרַכְתֶּם גַּם אֹתִי׃}
}
{
	\eRJE{
		Take your sheep also, your oxen also, as you spoke, and
		go. And you’ll bless me as well!”
	}
}
{
}

\Verse{33}
{
	\hRJE{וַתֶּחֱזַק מִצְרַיִם עַל הָעָם לְמַהֵר לְשַׁלְּחָם מִן הָאָרֶץ כִּי אָמְרוּ כֻּלָּנוּ מֵתִים׃}
}
{
	\eRJE{
		And Egypt was forceful on the people to hurry to let them
		go from the land, because they said, “We’re all dead!”
	}
}
{
}

\Verse{34}
{
	\hRJE{וַיִּשָּׂא הָעָם אֶת בְּצֵקוֹ טֶרֶם יֶחְמָץ מִשְׁאֲרֹתָם צְרֻרֹת בְּשִׂמְלֹתָם עַל שִׁכְמָם׃}
}
{
	\eRJE{
		And the people carried off its dough before it leavened,
		their bowls being wrapped in their garments on their
		shoulder.
	}
}
{
%	\footnote{
%		before it leavened. My colleague William Propp asks: Is the bread on the
%		night of the exodus leavened or unleavened? He answers that it is
%		leavened. The reason it does not rise is that the people carry off the
%		dough before it has time to rise. The ban on eating leaven on Passover
%		forever after this is a commemoration of this occurrence: “You shall eat
%		unleavened bread … because you went out from the land of Egypt in haste”
%		(Deut 16:3). He may be right, or perhaps we should understand this verse
%		to mean that the people carry off the dough before they have time to put
%		in the leaven, or that they deliberately do not leaven the dough because
%		they anticipate leaving in a hurry. In any case, the historical origin
%		of not eating leaven on Passover has associations beyond this
%		etiological story. See the comment on Exod 29:2.
%	}
}

\Verse{35}
{
	\hRJE{וּבְנֵי יִשְׂרָאֵל עָשׂוּ כִּדְבַר מֹשֶׁה וַיִּשְׁאֲלוּ מִמִּצְרַיִם כְּלֵי כֶסֶף וּכְלֵי זָהָב וּשְׂמָלֹת׃}
}
{
	\eRJE{
		And the children of Israel had done according to Moses’
		word, and they asked items of silver and items of gold and
		garments from Egypt.
	}
}
{
}

\Verse{36}
{
	\hRJE{וַיהֹוָה נָתַן אֶת חֵן הָעָם בְּעֵינֵי מִצְרַיִם וַיַּשְׁאִלוּם וַיְנַצְּלוּ אֶת מִצְרָיִם׃}
}
{
	\eRJE{
		And YHWH had put the people’s favor in the Egyptians’
		eyes, and they lent to them, and they despoiled Egypt.
	}
}
{
}

\Verse{37}
{
	\hP{וַיִּסְעוּ בְנֵי יִשְׂרָאֵל מֵרַעְמְסֵס סֻכֹּתָה כְּשֵׁשׁ מֵאוֹת אֶלֶף רַגְלִי הַגְּבָרִים לְבַד מִטָּף׃}
}
{
	\eP{
		And the children of Israel traveled from Rameses to
		Succoth, about six hundred thousand on foot—men, apart
		from infants.
	}
}
{
}

\Verse{38}
{
	\hP{וְגַם עֵרֶב רַב עָלָה אִתָּם וְצֹאן וּבָקָר מִקְנֶה כָּבֵד מְאֹד׃}
}
{
	\eP{
		And also a great mixture had gone up with them, and sheep
		and oxen, a very heavy livestock.
	}
}
{
%	\footnote{
%		heavy. The word that described the Pharaoh’s oppression and the force of
%		the plagues now recurs to describe, for the first time, something good:
%		the substantial quantity of possessions that the people are able to take
%		with them. It is as if this good is in proportion to the bad that they
%		have experienced.
%	}
}

\Verse{39}
{
	\hP{וַיֹּאפוּ אֶת הַבָּצֵק אֲשֶׁר הוֹצִיאוּ מִמִּצְרַיִם עֻגֹת מַצּוֹת כִּי לֹא חָמֵץ כִּי גֹרְשׁוּ מִמִּצְרַיִם וְלֹא יָכְלוּ לְהִתְמַהְמֵהַּ וְגַם צֵדָה לֹא עָשׂוּ לָהֶם׃}
}
{
	\eP{
		And they baked the dough that they brought out of Egypt:
		cakes of unleavened bread, because it had not leavened,
		because they were driven from Egypt and were not able to
		delay, and they also had not made provisions for
		themselves.
	}
}
{
}

\Verse{40}
{
	\hP{וּמוֹשַׁב בְּנֵי יִשְׂרָאֵל אֲשֶׁר יָשְׁבוּ בְּמִצְרָיִם שְׁלֹשִׁים שָׁנָה וְאַרְבַּע מֵאוֹת שָׁנָה׃}
}
{
	\eP{
		And the duration of the children of Israel that they lived
		in Egypt was thirty years and four hundred years.
	}
}
{
}

\Verse{41}
{
	\hP{וַיְהִי מִקֵּץ שְׁלֹשִׁים שָׁנָה וְאַרְבַּע מֵאוֹת שָׁנָה וַיְהִי בְּעֶצֶם הַיּוֹם הַזֶּה יָצְאוּ כּל צִבְאוֹת יְהֹוָה מֵאֶרֶץ מִצְרָיִם׃}
}
{
	\eP{
		And it was at the end of thirty years and four hundred
		years, and it was in that very day: all of YHWH’s masses
		went out from the land of Egypt.
	}
}
{
}

\Verse{42}
{
	\hP{לֵיל שִׁמֻּרִים הוּא לַיהֹוָה לְהוֹצִיאָם מֵאֶרֶץ מִצְרָיִם הוּא הַלַּיְלָה הַזֶּה לַיהֹוָה שִׁמֻּרִים לְכל בְּנֵי יִשְׂרָאֵל לְדֹרֹתָם׃}
}
{
	\eP{
		It is a night to be observed for YHWH for bringing them
		out from the land of Egypt. It is this night, to be
		observed for YHWH for all the children of Israel through
		their generations.
	}
}
{
}

\Verse{43}
{
	\hP{וַיֹּאמֶר יְהֹוָה אֶל מֹשֶׁה וְאַהֲרֹן זֹאת חֻקַּת הַפָּסַח כּל בֶּן נֵכָר לֹא יֹאכַל בּוֹ׃}
}
{
	\eP{
		And YHWH said to Moses and Aaron, “This is the law of the
		Passover: any foreigner shall not eat from it.
	}
}
{
}

\Verse{44}
{
	\hP{וְכל עֶבֶד אִישׁ מִקְנַת כָּסֶף וּמַלְתָּה אֹתוֹ אָז יֹאכַל בּוֹ׃}
}
{
	\eP{
		And every slave of a man, purchased with money: you shall
		circumcise him; then he shall eat from it.
	}
}
{
%	\footnote{
%		slave. It is extraordinary that, just two verses after reporting that
%		Israel went free from Egypt, there is a reference to the possibility of
%		an Israelite having a slave! The Torah does not forbid slavery. Nor does
%		it criticize the Egyptians for having slaves, but rather for maltreating
%		their slaves. The Torah rather initiates a process that eventually is to
%		bring about an end to all slavery on earth, for Israelites and for
%		everyone else. See the comments on Lev Footnote 25:43.
%	}
}

\Verse{45}
{
	\hP{תּוֹשָׁב וְשָׂכִיר לֹא יֹאכַל בּוֹ׃}
}
{
	\eP{A visitor and an employee shall not eat from it.}
}
{
%	\footnote{
%		A visitor and an employee shall not eat from it. These are
%		non-Israelites, who do not partake of the Passover sacrificial meal.
%		This does not apply to the Seder meal, which has been celebrated in
%		place of the sacrificial meal ever since the destruction of the Temple
%		in 70 C.E. brought an end to all sacrifices in Judaism.
%	}
}

\Verse{46}
{
	\hP{בְּבַיִת אֶחָד יֵאָכֵל לֹא תוֹצִיא מִן הַבַּיִת מִן הַבָּשָׂר חוּצָה וְעֶצֶם לֹא תִשְׁבְּרוּ בוֹ׃}
}
{
	\eP{
		It shall be eaten in one house; you shall not take any of
		the meat from the house outside. And you shall not break a
		bone from it.
	}
}
{
}

\Verse{47}
{
	\hP{כּל עֲדַת יִשְׂרָאֵל יַעֲשׂוּ אֹתוֹ׃}
}
{
	\eP{All of the congregation of Israel shall do it.}
}
{
}

\Verse{48}
{
	\hP{וְכִי יָגוּר אִתְּךָ גֵּר וְעָשָׂה פֶסַח לַיהֹוָה הִמּוֹל לוֹ כל זָכָר וְאָז יִקְרַב לַעֲשֹׂתוֹ וְהָיָה כְּאֶזְרַח הָאָרֶץ וְכל עָרֵל לֹא יֹאכַל בּוֹ׃}
}
{
	\eP{
		And if an alien will reside with you and will make a
		Passover to YHWH, let him be circumcised, every male, and
		then he may come forward to do it, and he will be like a
		citizen of the land, but everyone who is uncircumcised
		shall not eat from it.
	}
}
{
}

\Verse{49}
{
	\hP{תּוֹרָה אַחַת יִהְיֶה לָאֶזְרָח וְלַגֵּר הַגָּר בְּתוֹכְכֶם׃}
}
{
	\eP{
		There shall be one instruction for the citizen and for the
		alien who resides among you.”
	}
}
{
%	\footnote{
%		one instruction for the citizen and for the alien. “Instruction” is the
%		translation of Hebrew Tôrh because it denotes both “teaching” and “law.”
%		This is the first occurrence of the word “Torah” in the Torah. It is
%		impressive beyond description that its first appearance is in a verse
%		declaring that a foreigner residing among the people of Israel has the
%		same legal status as an Israelite. The context here concerns the
%		Passover statute, but this principle of treating a resident alien the
%		same as any citizen will be repeated many (about fifteen) times in the
%		Torah.
%	}
}

\Verse{50}
{
	\hP{וַיַּעֲשׂוּ כּל בְּנֵי יִשְׂרָאֵל כַּאֲשֶׁר צִוָּה יְהֹוָה אֶת מֹשֶׁה וְאֶת אַהֲרֹן כֵּן עָשׂוּ׃}
}
{
	\eP{
		And all the children of Israel did as YHWH had commanded
		Moses and Aaron. They did so.
	}
}
{
}

\Verse{51}
{
	\hP{וַיְהִי בְּעֶצֶם הַיּוֹם הַזֶּה הוֹצִיא יְהֹוָה אֶת בְּנֵי יִשְׂרָאֵל מֵאֶרֶץ מִצְרַיִם עַל צִבְאֹתָם׃}
}
{
	\eP{
		And it was in this very day: YHWH brought out the children
		of Israel from the land of Egypt by their masses.
	}
}
{
}
